%% ============================================================
%% appendix_refactor2.tex — 부록 C: 2차 Refactoring 계획 (전문)
%% 원문: Document/REFACTORING_PLAN_2.md
%% ============================================================
\chaptersummary{시뮬·피코·웹 세 트랙을 대상으로 한 2차 Refactoring 고도화 계획이다. 각 트랙의 개선 목표와 구조 개편 방향을 구체화한다.}
\chapter{2차 Refactoring 계획 (전문)}
\label{app:refactor2}

\textit{작성일 2026-06-26 --- 1차 계획(\ref{app:refactor1})의 연장선.}
1차가 ``비대한 파일$\cdot$함수를 책임 단위로 쪼개는'' 동작 보존 구조 분할이었다면,
2차는 그 위에서 \textbf{세 트랙 --- 시뮬레이터 $\cdot$ 피코+블록코딩 $\cdot$ 웹
인터페이스 ---}\,를 함께 끌어올린다. 구조 분할(1차 계승)에 더해 개발 워크플로$\cdot$학습자
경험 개선을 함께 다룬다. 원문: \code{Document/REFACTORING\_PLAN\_2.md}.

\section{기반 전략 (요약)}
\begin{enumerate}
  \item \textbf{시뮬레이터} --- (1.1) 요소별 코드 분할, (1.2) 단위 메쉬 로딩 후 개발자
  도구로 Transform 조작.
  \item \textbf{피코 + 블록코딩} --- (2.1) 요소 피코 코드 점검, (2.2) 미션별 블록 노출
  시스템 개선, (2.3) 미션 주제별 필요 블록 설계.
  \item \textbf{웹 인터페이스} --- (3.1) 반응형 웹, (3.2) 모바일 환경 최적화, (3.3) 디자인
  현대화.
\end{enumerate}

\section{트랙 1 --- 시뮬레이터 (\code{Web/simulation.js}, \code{Web/Mesh/})}

\subsection{1.1 요소별 코드 분할}
\code{simulation.js}(약 2,333줄)의 거대 클로저(\code{buildSim}/\code{setupSimulation})를
\textbf{요소 단위}(로버 부품$\cdot$발사대$\cdot$신호등$\cdot$LED$\cdot$총$\cdot$로켓$\cdot$음파$\cdot$OLED 등)로 분할한다. 1차 계획 유형 A의 \code{sim/} 서브시스템
(\code{context/assets/leds/movement/gun/rocket/traffic/waves/oled/render/audio/dispatch/topics})
구조를 실제로 도입해, 요소 하나를 독립적으로 보고 고칠 수 있게 한다. 이 작업은
1차 로드맵 6단계와 동일하며, 2차에서 본격 착수 항목으로 승격한다.

\subsection{1.2 단위 메쉬 로딩 후 Transform 조작}
부품 메쉬(\code{Web/Mesh/RoverParts/*.glb}, \code{EnvAssets/*} 등)를 \textbf{개별 로딩}한
뒤, 브라우저 개발자 도구로 위치$\cdot$회전$\cdot$스케일(Transform)을 실시간 조작해 정렬값을
잡고, 확정값을 코드(상수/설정)로 역반영하는 워크플로를 만든다~\cite{threejs, gltf, meshopt}.
하드코딩된 배치 좌표를 눈대중으로 맞추던 수작업을 도구화한다.
\begin{itemize}
  \item \textbf{구현 아이디어} --- 디버그 모드(\code{?debug=1})에서 three.js
  \code{TransformControls} 부착 → 선택 객체의 \code{position/rotation/scale}을 패널$\cdot$콘솔에
  출력 → 코드로 붙여넣기.
  \item \textbf{단위 메쉬 로더} --- 부품을 하나씩 로딩$\cdot$점검하는 경로로 모델 교체$\cdot$정렬
  검증을 빠르게.
\end{itemize}

\section{트랙 2 --- 피코 + 블록코딩 (\code{Pico/}, \code{blocklyconfig.js}, \code{Missions/})}

\subsection{2.1 요소 피코 코드 점검}
\code{process\_data.py} 핸들러와 하드웨어 모듈(모터$\cdot$LED$\cdot$부저$\cdot$거리$\cdot$자기$\cdot$총$\cdot$OLED)을 요소별로 점검$\cdot$정리한다(1차 유형 B 분할과 연계). 각 명령
문자열이 실제 하드웨어 동작과 1:1로 맞는지 검증하고, 죽은 코드$\cdot$중복 분기를
솎아낸다.

\subsection{2.2 미션별 블록 노출 시스템 개선}
현재 toolbox는 모든 블록을 항상 노출한다(\code{main.html} 정적 XML). 미션(주제)에 따라
\textbf{필요한 블록만 보여주는 노출 제어 시스템}을 도입한다. 정적 toolbox 구성을
\textbf{미션 메타데이터 기반 동적 구성}으로 전환한다(\code{Missions/}, XML Presets,
\code{MISSIONS.md} 연계).

\subsection{2.3 미션 주제별 필요 블록 설계}
미션 맥락에 맞는 블록 세트를 설계한다. 같은 ``전진''이라도 맥락에 따라 구분한다.

\begin{center}
\renewcommand{\arraystretch}{1.25}
\begin{tabularx}{\linewidth}{@{}>{\sffamily}L{0.24\linewidth} >{\sffamily}L{0.32\linewidth} L{0.3\linewidth}@{}}
\toprule
\sffamily 맥락 & \sffamily 블록 & \sffamily 핵심 파라미터 \\
\midrule
알비 DC 전진 & DC 모터 주행 & 방향(전/후) \\
속도 조절 & 주행 속도 설정 & speed 값 \\
발사 동작 & 로켓 발사 시퀀스 & 카운트$\cdot$점화 \\
\bottomrule
\end{tabularx}
\end{center}

\noindent 맥락에 따라 블록의 명칭$\cdot$파라미터$\cdot$노출을 다르게 설계해, 학습자가 미션
목표에 집중하도록 한다.

\section{트랙 3 --- 웹 인터페이스 (\code{Web/*.html}, \code{index.css}, \code{mobile-preview.js})}
\begin{itemize}
  \item \textbf{3.1 반응형 웹} --- 데스크톱$\cdot$태블릿$\cdot$모바일 폭에 맞춰 레이아웃을
  재구성. 현재 \code{index.css}는 \code{@media} 일부만 대응 → 체계적 브레이크포인트로 정리.
  \item \textbf{3.2 모바일 최적화} --- 모바일에서 핵심 동작에 필요한 최소 정보만 노출.
  \code{mobile-preview.js}의 \code{?mobile=true} 미리보기를 실제 모바일 UX로 발전(불필요한
  패널 축소, 터치 타깃 확대)~\cite{webbluetooth}.
  \item \textbf{3.3 디자인 현대화} --- 타이포$\cdot$여백$\cdot$컬러$\cdot$컴포넌트를 현대적 디자인
  시스템으로 정리, 홈 3D 히어로(우주 유영 연출)와 톤 통일.
\end{itemize}

\section{우선순위 로드맵}
\begin{center}
\renewcommand{\arraystretch}{1.25}
\begin{tabularx}{\linewidth}{@{}c >{\sffamily}l L{0.44\linewidth} c >{\footnotesize}L{0.16\linewidth}@{}}
\toprule
\sffamily 순서 & \sffamily 트랙 & \sffamily 작업 & \sffamily 위험 & \sffamily 효과 \\
\midrule
1 & 웹 & 3.1 반응형 $\cdot$ 3.3 디자인 현대화 & 낮음 & 가시 효과 즉시 \\
2 & 웹 & 3.2 모바일 최소 정보 추출 & 중 & 모바일 사용성↑ \\
3 & 블록 & 2.2 미션별 블록 노출 & 중 & 학습 집중도↑ \\
4 & 피코 & 2.1 요소 피코 코드 점검 & 중 & 동작 신뢰성↑ \\
5 & 시뮬 & 1.1 요소별 분할 $\cdot$ 1.2 Transform 도구 & 높음 & 최대 가치 \\
6 & 블록 & 2.3 미션 주제별 블록 설계 & 중 & 커리큘럼 정합성↑ \\
\bottomrule
\end{tabularx}
\end{center}

\noindent\textbf{원칙:} 동작 보존 분할(1차 계승)과 사용자 경험 개선은 \textbf{커밋을
분리}한다. 가시 효과$\cdot$저위험(웹)부터 시작해, 가장 어려운 시뮬레이터(1차 6단계)는
패턴 검증 후 착수한다.
